\section{背景知识：物理内存布局}

BIOS 通过 Multiboot2 Memory Map Tag（type=6）报告物理内存区域。
一台典型的 256MiB QEMU 虚拟机的内存映射如下：

\begin{figure}[H]
  \centering
  % figures/ch07/fig01-256mib-qemu-vm.tex — auto-extracted from chapter source
\begin{tikzpicture}[node distance=0cm]
  % 地址轴
  \draw[-{Stealth}, thick] (0, 0) -- (12, 0) node[right, font=\small] {物理地址};

  % 区域
  \fill[green!20] (0, 0.2) rectangle (1.2, 0.8);
  \node[font=\tiny] at (0.6, 0.5) {640K};
  \node[font=\tiny, below] at (0.6, 0.15) {\hex{0}};

  \fill[red!20] (1.2, 0.2) rectangle (2, 0.8);
  \node[font=\tiny] at (1.6, 0.5) {保留};
  \node[font=\tiny, below] at (1.6, 0.15) {VGA/ROM};

  \fill[green!30] (2, 0.2) rectangle (10, 0.8);
  \node[font=\small] at (6, 0.5) {可用 RAM（\textasciitilde 126 MiB）};
  \node[font=\tiny, below] at (2, 0.15) {\hex{100000}};
  \node[font=\tiny, below] at (10, 0.15) {\hex{7FE0000}};

  \fill[orange!20] (10, 0.2) rectangle (12, 0.8);
  \node[font=\tiny] at (11, 0.5) {ACPI/保留};
\end{tikzpicture}

  \caption{典型物理内存布局（256MiB QEMU VM）}
\end{figure}

\subsection{Multiboot2 Memory Map Tag 格式}

\codefile{src/kernel/mm/mmap.c}
\begin{lstlisting}[style=cstyle]
struct mb2_tag_mmap {
    uint32_t type;          // = 6 (MB2_TAG_TYPE_MMAP)
    uint32_t size;          // tag 总字节数
    uint32_t entry_size;    // 每个 entry 的字节数（通常 24）
    uint32_t entry_version; // = 0
    // 后面紧跟 N 个 mb2_mmap_entry
};

struct mb2_mmap_entry {
    uint64_t addr;          // 区域起始物理地址
    uint64_t len;           // 区域长度（字节）
    uint32_t type;          // 1=Available, 2=Reserved, 3=ACPI, ...
    uint32_t reserved;
};
\end{lstlisting}

\begin{thinkbox}
\texttt{addr} 和 \texttt{len} 是 64 位的。在 32 位内核中，我们只能访问低 4GiB
的物理内存。如何处理报告了 4GiB 以上 RAM 的机器？
（答：截断到 \hex{FFFFFFFF}——安全忽略高于 4GiB 的区域。）
\end{thinkbox}

% ============================================================================

